\section*{Capítulo \ref{ch:g5:irreversible-changes} --- Transformações que não podem ser desfeitas}

\begin{solution}{exo:g5:irreversible-changes:1}
Derreter chocolate: pode ser desfeito (ele endurece de novo quando
esfria). Queimar um fósforo: não pode. \omterm{def:g2:mixing-and-dissolving:dissolve}{Dissolver} sal: pode ser desfeito
(fazendo a água evaporar). Torrar pão: não pode.
\end{solution}

\begin{solution}{exo:g5:irreversible-changes:2}
Uma transformação em que aparecem substâncias novas. Na cozinha:
cozinhar um ovo, assar um bolo (e também: torrar pão, uma maçã cortada
que escurece).
\end{solution}

\begin{solution}{exo:g5:irreversible-changes:3}
Fazendo a água evaporar: o sal fica no recipiente. \omterm{def:g2:mixing-and-dissolving:dissolve}{Dissolver} é uma
\omterm{def:g5:irreversible-changes:reversible}{transformação reversível}.
\end{solution}

\begin{solution}{exo:g5:irreversible-changes:4}
Um cheiro novo, e grumos (um sólido) aparecendo no leite.
\end{solution}

\begin{solution}{exo:g5:irreversible-changes:5}
As transformações da esquerda: o açúcar se dissolvendo na água, o gelo
derretendo, a areia \omterm{def:g2:mixing-and-dissolving:mix}{misturada} com o cascalho. A seta dupla quer dizer
que a transformação pode ir nos dois sentidos: ela pode ser desfeita.
\end{solution}

\begin{solution}{exo:g5:irreversible-changes:6}
O óleo mantém o \omterm{def:g4:air-a-mixture-of-gases:air}{ar} e a água longe do ferro da corrente: ele deixa a
ferrugem mais lenta.
\end{solution}

\begin{solution}{exo:g5:irreversible-changes:7}
Bolhas de gás aparecendo num líquido sem nenhum aquecimento.
\end{solution}

\begin{solution}{exo:g5:irreversible-changes:8}
Úteis: cozinhar os alimentos, assar o pão, queimar gás para aquecer.
Prejudiciais: a ferrugem, a comida que estraga, a madeira que apodrece.
\end{solution}

\begin{solution}{exo:g5:irreversible-changes:9}
O frio deixa mais lentas as \omterm{def:g5:irreversible-changes:chemical-change}{transformações químicas} que fazem o leite
azedar, e assim ele dura mais.
\end{solution}

\begin{solution}{exo:g5:irreversible-changes:10}
Não. As bolhas são vapor de água, que volta a ser água quando esfria:
nada de novo é fabricado, e a transformação pode ser desfeita. As bolhas
sozinhas não provam que apareceu uma substância nova.
\end{solution}

\begin{solution}{exo:g5:irreversible-changes:11}
Uma vela acesa libera calor e luz, e fabrica substâncias novas: vapor de
água (um copo frio fica embaçado acima dela) e dióxido de carbono. A
cera que \omterm{def:g4:what-a-fire-needs:burning}{queima} não volta. Já a cera que derrete endurece de novo quando
esfria.
\end{solution}
